%! AP Statistics — Unit 3, Lesson 3.2 — Homework
\documentclass[10pt]{article}
\usepackage{apstats-article}
\usepackage{apstats-boxes}
\usepackage{grading}
\gradingassignment{apstats-u03-l02-hw}
\gradingcourse{advanced_statistics}

\begin{document}

\pageheader{Unit 3, Lesson 3.2}{Homework}
{\small\itshape Write each final answer in its box. Only what is inside a box is graded.}\par

\vspace{0.1in}

\begin{practicebox}
\small
For each study, classify it as an observational study or an experiment, identify
whether random selection and/or random assignment were used, and state what
conclusions are valid.

\vspace{8pt}
\textbf{1.} A nutritionist randomly selects 500 adults from the city directory and
asks each one about fruit and vegetable consumption and blood pressure. High consumers
have significantly lower blood pressure.

\vspace{6pt}
Study type: \makebox[5cm][r]{\smash{\gradedblank[id=hw1a, type=short, points=1, answer={observational study}, check=text]{1.3in}{0.2in}}} \quad Random selection? \makebox[2cm][l]{\hspace{31bp}\smash{\gradedblank[id=hw1b, type=short, points=1, answer=yes, check=text]{0.45in}{0.2in}}}\\[6pt]
Random assignment? \makebox[2cm][l]{\hspace{0.6in}\smash{\gradedblank[id=hw1c, type=short, points=1, answer=no, check=text]{0.8in}{0.2in}}}\\[6pt]
Can we conclude diet \emph{causes} lower BP? \makebox[3cm][l]{\hspace{44bp}\smash{\gradedblank[id=hw1d, type=short, points=1, answer=no, check=text]{0.65in}{0.2in}}} \quad Generalize to city adults? \makebox[3cm][l]{\hspace{46bp}\smash{\gradedblank[id=hw1e, type=short, points=1, answer=yes, check=text]{0.9in}{0.2in}}}\\[6pt]
Explain:\par
\gradedhide{\gradedhere{\gradedblank[id=hw1f, type=rubric, points=2, display={Observational: confounding (exercise, genetics), so no causation. Random selection: generalize to city adults.}]{\linewidth+6pt}{2in}}\writelines{9}}

\vspace{8pt}
\textbf{2.} A tech company recruits 80 employees and randomly assigns 40 to a
standing-desk program and 40 to regular desks. After 3 months, the standing-desk
group reports significantly less back pain.

\vspace{6pt}
Study type: \makebox[5cm][r]{\smash{\gradedblank[id=hw2a, type=short, points=1, answer=experiment, check=text]{1.3in}{0.2in}}} \quad Random selection? \makebox[2cm][l]{\hspace{31bp}\smash{\gradedblank[id=hw2b, type=short, points=1, answer=no, check=text]{0.45in}{0.2in}}}\\[6pt]
Random assignment? \makebox[2cm][l]{\hspace{0.6in}\smash{\gradedblank[id=hw2c, type=short, points=1, answer=yes, check=text]{0.8in}{0.2in}}}\\[6pt]
Can we conclude standing desks \emph{caused} less pain? \makebox[3cm][l]{\hspace{46bp}\smash{\gradedblank[id=hw2d, type=short, points=1, answer=yes, check=text]{1in}{0.2in}}}\\[6pt]
Can we generalize to all office workers? \makebox[3cm][l]{\hspace{46bp}\smash{\gradedblank[id=hw2e, type=short, points=1, answer=no, check=text]{1in}{0.2in}}}\\[6pt]
Explain:\par
\gradedhide{\gradedhere{\gradedblank[id=hw2f, type=rubric, points=2, display={Random assignment, so causation. Volunteers, not randomly selected: no generalizing beyond this company.}]{\linewidth+6pt}{2in}}\writelines{9}}

\tcbbreak
\vspace{8pt}
\textbf{3.} Free Response Practice:\\
\textit{A researcher wants to determine whether a new reading intervention improves
comprehension. She has 60 third-grade students from a single school who volunteer
to participate. She randomly assigns 30 to the intervention and 30 to a control
group. After 6 weeks, students in the intervention group score significantly higher.}

\vspace{6pt}
\begin{enumerate}[label=(\alph*), leftmargin=2em, itemsep=8pt]
  \item Can we conclude the intervention caused higher scores? Explain.\\
        \gradedhide{\gradedhere{\gradedblank[id=hw3a, type=rubric, points=2, display={Yes: random assignment, so the intervention caused the higher scores.}]{\linewidth+6pt}{2in}}\writeline\\[1pt]\writelines{9}}
  \item Can we generalize to all third-graders? Explain.\\
        \gradedhide{\gradedhere{\gradedblank[id=hw3b, type=rubric, points=2, display={No: volunteers from one school, not a random sample of all third-graders.}]{\linewidth+6pt}{2in}}\writeline\\[1pt]\writelines{9}}
  \item What change to the study design would allow generalization?\\
        \gradedhide{\gradedhere{\gradedblank[id=hw3c, type=rubric, points=2, display={Randomly select students from all third-graders (e.g., a district registry).}]{\linewidth+6pt}{2in}}\writeline\\[1pt]\writelines{9}}
\end{enumerate}
\end{practicebox}

\end{document}
